% =====================================================================
\chapter{Introdução às equações diferenciais}

Muitas leis da natureza não dizem diretamente quanto vale uma grandeza, mas \emph{como ela
varia}: a velocidade de queda depende da velocidade atual, uma população cresce em proporção
ao seu tamanho, um corpo esfria a uma taxa proporcional à diferença de temperatura para o
ambiente. Equações que envolvem uma função desconhecida e suas derivadas são as
\emph{equações diferenciais}. Neste capítulo resolvemos as mais simples (separáveis, logística
e lineares de primeira ordem) e vemos como estudá-las mesmo sem fórmula, pelo campo de
direções e pelo método de Euler.

% ---------------------------------------------------------------------
\section{Fundamentos das equações diferenciais}

\begin{definicao}
Uma \textbf{equação diferencial}\index{equacao diferencial@equação diferencial} é uma equação
que envolve uma função desconhecida $y=f(x)$ e uma ou mais de suas derivadas. A
\textbf{ordem}\index{ordem de uma equação diferencial} da equação é a ordem da maior derivada
que nela aparece. Uma \textbf{solução} é uma função $y=f(x)$ que satisfaz a equação para todo
$x$ de algum intervalo.
\end{definicao}

Por exemplo, $y'=3x^2$ tem ordem 1, $y''+4y=0$ tem ordem 2 e $(y''')^2-4y'=x$ tem ordem 3
(o expoente 2 não muda a ordem). Para verificar uma solução candidata, basta derivá-la e
substituir.

\begin{exemplo}{1.1}{Verificando uma solução}
Mostre que $y=e^{-3x}+2x+3$ é solução de $y'+3y=6x+11$.
\solucao
$y'=-3e^{-3x}+2$, logo $y'+3y=-3e^{-3x}+2+3e^{-3x}+6x+9=6x+11$. \checkmark
\end{exemplo}

\textbf{Solução geral e solução particular.} A equação $y'=2x$ tem como soluções todas as
funções $y=x^2+C$: uma \emph{família} a um parâmetro, a \textbf{solução geral}\index{solucao
geral@solução geral}. Cada escolha de $C$ dá uma \textbf{solução particular}. Em geral, uma
equação de ordem $n$ tem solução geral com $n$ constantes arbitrárias.

\begin{definicao}[Problema de valor inicial]
Um \textbf{problema de valor inicial} (PVI)\index{problema de valor inicial} é uma equação
diferencial acompanhada de condições que fixam as constantes, como $y(x_0)=y_0$ (e, para
ordem 2, também $y'(x_0)=y_1$).
\end{definicao}

\begin{exemplo}{1.2}{Resolvendo um PVI por integração direta}
(a) $y'=3e^x+x^2-4$, $y(0)=5$. Integrando, $y=3e^x+\frac{x^3}3-4x+C$; de $y(0)=3+C=5$
vem $C=2$: $y=3e^x+\frac{x^3}3-4x+2$.

(b) Uma bola é lançada para cima a $10$ m/s de uma altura de $3$ m. Com $a=-9{,}8$ m/s$^2$:
$v(t)=-9{,}8t+10$ e $s(t)=-4{,}9t^2+10t+3$. A altura máxima ocorre em $v=0$, isto é,
$t\approx1{,}02$ s, quando $s\approx8{,}10$ m.
\end{exemplo}

% ---------------------------------------------------------------------
\section{Campos de direções e o método de Euler}

\subsection{Campos de direções}

Uma equação de primeira ordem $y'=f(x,y)$ informa, em cada ponto $(x,y)$ do plano, a
\emph{inclinação} que uma solução que passa por ali deve ter. Desenhando pequenos segmentos com
essas inclinações numa grade obtemos o \textbf{campo de direções}\index{campo de direções}.
As soluções são as curvas que seguem o campo, tangentes aos segmentos em todos os pontos.

\figtikz{\begin{tikzpicture}
\begin{axis}[painel,width=8.4cm,height=6.4cm,xmin=-2.2,xmax=3.2,ymin=-2.2,ymax=3.2,xtick={-2,-1,1,2,3},ytick={-2,-1,1,2,3}]
\addplot[cinza,thin,quiver={u=\thisrow{u},v=\thisrow{v}}] table[row sep=\\] {x y u v\\
-2.119 -2.000 0.2370 0.0000\\
-2.101 -1.462 0.2020 -0.1674\\
-2.075 -0.967 0.1497 -0.2481\\
-2.057 -0.496 0.1131 -0.2812\\
-2.045 -0.034 0.0894 -0.2963\\
-2.037 0.425 0.0735 -0.3042\\
-2.031 0.882 0.0621 -0.3088\\
-2.027 1.338 0.0537 -0.3117\\
-2.024 1.793 0.0473 -0.3136\\
-2.021 2.248 0.0422 -0.3149\\
-2.019 2.703 0.0381 -0.3158\\
-2.017 3.158 0.0347 -0.3165\\
-1.646 -2.084 0.2020 0.1674\\
-1.664 -1.545 0.2370 0.0000\\
-1.646 -1.007 0.2020 -0.1674\\
-1.620 -0.512 0.1497 -0.2481\\
-1.602 -0.041 0.1131 -0.2812\\
-1.590 0.421 0.0894 -0.2963\\
-1.582 0.879 0.0735 -0.3042\\
-1.577 1.336 0.0621 -0.3088\\
-1.572 1.792 0.0537 -0.3117\\
-1.569 2.248 0.0473 -0.3136\\
-1.567 2.703 0.0422 -0.3149\\
-1.565 3.158 0.0381 -0.3158\\
-1.166 -2.124 0.1497 0.2481\\
-1.192 -1.629 0.2020 0.1674\\
-1.209 -1.091 0.2370 0.0000\\
-1.192 -0.553 0.2020 -0.1674\\
-1.166 -0.058 0.1497 -0.2481\\
-1.147 0.413 0.1131 -0.2812\\
-1.136 0.875 0.0894 -0.2963\\
-1.128 1.334 0.0735 -0.3042\\
-1.122 1.791 0.0621 -0.3088\\
-1.118 2.247 0.0537 -0.3117\\
-1.115 2.702 0.0473 -0.3136\\
-1.112 3.157 0.0422 -0.3149\\
-0.693 -2.141 0.1131 0.2812\\
-0.711 -1.669 0.1497 0.2481\\
-0.737 -1.175 0.2020 0.1674\\
-0.755 -0.636 0.2370 0.0000\\
-0.737 -0.098 0.2020 -0.1674\\
-0.711 0.397 0.1497 -0.2481\\
-0.693 0.868 0.1131 -0.2812\\
-0.681 1.330 0.0894 -0.2963\\
-0.673 1.788 0.0735 -0.3042\\
-0.667 2.245 0.0621 -0.3088\\
-0.663 2.701 0.0537 -0.3117\\
-0.660 3.157 0.0473 -0.3136\\
-0.227 -2.148 0.0894 0.2963\\
-0.238 -1.686 0.1131 0.2812\\
-0.257 -1.215 0.1497 0.2481\\
-0.283 -0.720 0.2020 0.1674\\
-0.300 -0.182 0.2370 0.0000\\
-0.283 0.356 0.2020 -0.1674\\
-0.257 0.851 0.1497 -0.2481\\
-0.238 1.322 0.1131 -0.2812\\
-0.227 1.785 0.0894 -0.2963\\
-0.219 2.243 0.0735 -0.3042\\
-0.213 2.700 0.0621 -0.3088\\
-0.209 3.156 0.0537 -0.3117\\
0.236 -2.152 0.0735 0.3042\\
0.228 -1.694 0.0894 0.2963\\
0.216 -1.232 0.1131 0.2812\\
0.198 -0.760 0.1497 0.2481\\
0.172 -0.266 0.2020 0.1674\\
0.154 0.273 0.2370 0.0000\\
0.172 0.811 0.2020 -0.1674\\
0.198 1.306 0.1497 -0.2481\\
0.216 1.777 0.1131 -0.2812\\
0.228 2.239 0.0894 -0.2963\\
0.236 2.698 0.0735 -0.3042\\
0.242 3.154 0.0621 -0.3088\\
0.696 -2.154 0.0621 0.3088\\
0.691 -1.698 0.0735 0.3042\\
0.683 -1.239 0.0894 0.2963\\
0.671 -0.777 0.1131 0.2812\\
0.652 -0.306 0.1497 0.2481\\
0.626 0.189 0.2020 0.1674\\
0.609 0.727 0.2370 0.0000\\
0.626 1.266 0.2020 -0.1674\\
0.652 1.760 0.1497 -0.2481\\
0.671 2.232 0.1131 -0.2812\\
0.683 2.694 0.0894 -0.2963\\
0.691 3.152 0.0735 -0.3042\\
1.155 -2.156 0.0537 0.3117\\
1.151 -1.700 0.0621 0.3088\\
1.145 -1.243 0.0735 0.3042\\
1.137 -0.785 0.0894 0.2963\\
1.125 -0.322 0.1131 0.2812\\
1.107 0.149 0.1497 0.2481\\
1.081 0.644 0.2020 0.1674\\
1.063 1.182 0.2370 0.0000\\
1.081 1.720 0.2020 -0.1674\\
1.107 2.215 0.1497 -0.2481\\
1.125 2.686 0.1131 -0.2812\\
1.137 3.148 0.0894 -0.2963\\
1.613 -2.157 0.0473 0.3136\\
1.609 -1.701 0.0537 0.3117\\
1.605 -1.245 0.0621 0.3088\\
1.600 -0.788 0.0735 0.3042\\
1.592 -0.330 0.0894 0.2963\\
1.580 0.132 0.1131 0.2812\\
1.561 0.603 0.1497 0.2481\\
1.535 1.098 0.2020 0.1674\\
1.518 1.636 0.2370 0.0000\\
1.535 2.175 0.2020 -0.1674\\
1.561 2.669 0.1497 -0.2481\\
1.580 3.141 0.1131 -0.2812\\
2.070 -2.157 0.0422 0.3149\\
2.067 -1.702 0.0473 0.3136\\
2.064 -1.247 0.0537 0.3117\\
2.060 -0.791 0.0621 0.3088\\
2.054 -0.334 0.0735 0.3042\\
2.046 0.125 0.0894 0.2963\\
2.034 0.587 0.1131 0.2812\\
2.016 1.058 0.1497 0.2481\\
1.990 1.553 0.2020 0.1674\\
1.972 2.091 0.2370 0.0000\\
1.990 2.629 0.2020 -0.1674\\
2.016 3.124 0.1497 -0.2481\\
2.526 -2.158 0.0381 0.3158\\
2.524 -1.703 0.0422 0.3149\\
2.522 -1.248 0.0473 0.3136\\
2.519 -0.792 0.0537 0.3117\\
2.514 -0.336 0.0621 0.3088\\
2.509 0.121 0.0735 0.3042\\
2.501 0.579 0.0894 0.2963\\
2.489 1.041 0.1131 0.2812\\
2.471 1.512 0.1497 0.2481\\
2.444 2.007 0.2020 0.1674\\
2.427 2.545 0.2370 0.0000\\
2.444 3.084 0.2020 -0.1674\\
2.983 -2.158 0.0347 0.3165\\
2.981 -1.703 0.0381 0.3158\\
2.979 -1.248 0.0422 0.3149\\
2.976 -0.793 0.0473 0.3136\\
2.973 -0.338 0.0537 0.3117\\
2.969 0.118 0.0621 0.3088\\
2.963 0.575 0.0735 0.3042\\
2.955 1.034 0.0894 0.2963\\
2.943 1.496 0.1131 0.2812\\
2.925 1.967 0.1497 0.2481\\
2.899 2.462 0.2020 0.1674\\
2.881 3.000 0.2370 0.0000\\
};
\foreach \c in {-2,0,1,3}{\edef\tmp{\noexpand\addplot[laranja,thick,domain=-2.2:3.2,samples=80]{x-1+\c*exp(-x)};}\tmp}
\addplot[azul,very thick,domain=-2.2:3.2]{x-1};
\end{axis}\end{tikzpicture}}{Campo de direções de $y'=x-y$ e algumas soluções $y=x-1+Ce^{-x}$
(verde). Todas se aproximam da reta $y=x-1$ (grafite), que também é solução ($C=0$).}

\textbf{Soluções de equilíbrio.} Numa equação \emph{autônoma} $y'=f(y)$, cada zero $c$ de $f$
dá uma solução constante $y\equiv c$, uma \textbf{solução de equilíbrio}\index{solucao de
equilibrio@solução de equilíbrio}. Ela é \emph{estável} (assintoticamente) se as soluções
próximas se aproximam dela, \emph{instável} se se afastam e \emph{semiestável} se se aproximam
de um lado e se afastam do outro. Basta estudar o sinal de $f$.

\begin{exemplo}{1.3}{Classificando equilíbrios}
Para $y'=(y-2)(y+1)$, os equilíbrios são $y=-1$ e $y=2$. Para $y<-1$, $y'>0$ (as soluções
sobem até $-1$); entre $-1$ e $2$, $y'<0$ (descem até $-1$); para $y>2$, $y'>0$ (sobem
indefinidamente). Logo $y=-1$ é estável e $y=2$ é instável.
\end{exemplo}

\figtikz{\begin{tikzpicture}
\begin{axis}[painel,width=8.4cm,height=6cm,xmin=0,xmax=4,ymin=-3,ymax=3.5,xtick={1,2,3},ytick={-1,2},xlabel={$t$}]
\addplot[cinza,thin,quiver={u=\thisrow{u},v=\thisrow{v}}] table[row sep=\\] {x y u v\\
0.194 -2.950 0.0120 0.2993\\
0.190 -2.395 0.0195 0.2982\\
0.180 -1.837 0.0397 0.2923\\
0.129 -1.225 0.1427 0.1764\\
0.158 -0.450 0.0844 -0.2634\\
0.175 0.116 0.0504 -0.2875\\
0.178 0.672 0.0446 -0.2902\\
0.174 1.225 0.0519 -0.2867\\
0.154 1.764 0.0923 -0.2557\\
0.139 2.083 0.1226 0.2158\\
0.182 2.599 0.0364 0.2936\\
0.191 3.151 0.0185 0.2984\\
0.594 -2.950 0.0120 0.2993\\
0.590 -2.395 0.0195 0.2982\\
0.580 -1.837 0.0397 0.2923\\
0.529 -1.225 0.1427 0.1764\\
0.558 -0.450 0.0844 -0.2634\\
0.575 0.116 0.0504 -0.2875\\
0.578 0.672 0.0446 -0.2902\\
0.574 1.225 0.0519 -0.2867\\
0.554 1.764 0.0923 -0.2557\\
0.539 2.083 0.1226 0.2158\\
0.582 2.599 0.0364 0.2936\\
0.591 3.151 0.0185 0.2984\\
0.994 -2.950 0.0120 0.2993\\
0.990 -2.395 0.0195 0.2982\\
0.980 -1.837 0.0397 0.2923\\
0.929 -1.225 0.1427 0.1764\\
0.958 -0.450 0.0844 -0.2634\\
0.975 0.116 0.0504 -0.2875\\
0.978 0.672 0.0446 -0.2902\\
0.974 1.225 0.0519 -0.2867\\
0.954 1.764 0.0923 -0.2557\\
0.939 2.083 0.1226 0.2158\\
0.982 2.599 0.0364 0.2936\\
0.991 3.151 0.0185 0.2984\\
1.394 -2.950 0.0120 0.2993\\
1.390 -2.395 0.0195 0.2982\\
1.380 -1.837 0.0397 0.2923\\
1.329 -1.225 0.1427 0.1764\\
1.358 -0.450 0.0844 -0.2634\\
1.375 0.116 0.0504 -0.2875\\
1.378 0.672 0.0446 -0.2902\\
1.374 1.225 0.0519 -0.2867\\
1.354 1.764 0.0923 -0.2557\\
1.339 2.083 0.1226 0.2158\\
1.382 2.599 0.0364 0.2936\\
1.391 3.151 0.0185 0.2984\\
1.794 -2.950 0.0120 0.2993\\
1.790 -2.395 0.0195 0.2982\\
1.780 -1.837 0.0397 0.2923\\
1.729 -1.225 0.1427 0.1764\\
1.758 -0.450 0.0844 -0.2634\\
1.775 0.116 0.0504 -0.2875\\
1.778 0.672 0.0446 -0.2902\\
1.774 1.225 0.0519 -0.2867\\
1.754 1.764 0.0923 -0.2557\\
1.739 2.083 0.1226 0.2158\\
1.782 2.599 0.0364 0.2936\\
1.791 3.151 0.0185 0.2984\\
2.194 -2.950 0.0120 0.2993\\
2.190 -2.395 0.0195 0.2982\\
2.180 -1.837 0.0397 0.2923\\
2.129 -1.225 0.1427 0.1764\\
2.158 -0.450 0.0844 -0.2634\\
2.175 0.116 0.0504 -0.2875\\
2.178 0.672 0.0446 -0.2902\\
2.174 1.225 0.0519 -0.2867\\
2.154 1.764 0.0923 -0.2557\\
2.139 2.083 0.1226 0.2158\\
2.182 2.599 0.0364 0.2936\\
2.191 3.151 0.0185 0.2984\\
2.594 -2.950 0.0120 0.2993\\
2.590 -2.395 0.0195 0.2982\\
2.580 -1.837 0.0397 0.2923\\
2.529 -1.225 0.1427 0.1764\\
2.558 -0.450 0.0844 -0.2634\\
2.575 0.116 0.0504 -0.2875\\
2.578 0.672 0.0446 -0.2902\\
2.574 1.225 0.0519 -0.2867\\
2.554 1.764 0.0923 -0.2557\\
2.539 2.083 0.1226 0.2158\\
2.582 2.599 0.0364 0.2936\\
2.591 3.151 0.0185 0.2984\\
2.994 -2.950 0.0120 0.2993\\
2.990 -2.395 0.0195 0.2982\\
2.980 -1.837 0.0397 0.2923\\
2.929 -1.225 0.1427 0.1764\\
2.958 -0.450 0.0844 -0.2634\\
2.975 0.116 0.0504 -0.2875\\
2.978 0.672 0.0446 -0.2902\\
2.974 1.225 0.0519 -0.2867\\
2.954 1.764 0.0923 -0.2557\\
2.939 2.083 0.1226 0.2158\\
2.982 2.599 0.0364 0.2936\\
2.991 3.151 0.0185 0.2984\\
3.394 -2.950 0.0120 0.2993\\
3.390 -2.395 0.0195 0.2982\\
3.380 -1.837 0.0397 0.2923\\
3.329 -1.225 0.1427 0.1764\\
3.358 -0.450 0.0844 -0.2634\\
3.375 0.116 0.0504 -0.2875\\
3.378 0.672 0.0446 -0.2902\\
3.374 1.225 0.0519 -0.2867\\
3.354 1.764 0.0923 -0.2557\\
3.339 2.083 0.1226 0.2158\\
3.382 2.599 0.0364 0.2936\\
3.391 3.151 0.0185 0.2984\\
3.794 -2.950 0.0120 0.2993\\
3.790 -2.395 0.0195 0.2982\\
3.780 -1.837 0.0397 0.2923\\
3.729 -1.225 0.1427 0.1764\\
3.758 -0.450 0.0844 -0.2634\\
3.775 0.116 0.0504 -0.2875\\
3.778 0.672 0.0446 -0.2902\\
3.774 1.225 0.0519 -0.2867\\
3.754 1.764 0.0923 -0.2557\\
3.739 2.083 0.1226 0.2158\\
3.782 2.599 0.0364 0.2936\\
3.791 3.151 0.0185 0.2984\\
};
\addplot[azul,very thick,domain=0:4]{-1}; \addplot[azul,very thick,dashed,domain=0:4]{2};
\foreach \c in {-2.5,-1.6,0,1.5,1.9}{\pgfmathsetmacro{\A}{(\c-2)/(\c+1)}\edef\tmp{\noexpand\addplot[laranja,thick,domain=0:4,samples=120]{(2+\A*exp(3*x))/(1-\A*exp(3*x))};}\tmp}
\addplot[laranja,thick,domain=0:0.84,samples=120]{(2+0.0625*exp(3*x))/(1-0.0625*exp(3*x))};
\end{axis}\end{tikzpicture}}{Equilíbrios de $y'=(y-2)(y+1)$: $y=-1$ é estável (as soluções
vizinhas convergem para ela) e $y=2$ é instável (tracejada: as soluções se afastam).}

\subsection{O método de Euler}

Quando não há fórmula para a solução, aproximamos seguindo o campo de direções em pequenos
passos de tamanho $h$. Partindo de $(x_0,y_0)$, andamos sobre a reta tangente:

\begin{regra}{Método de Euler}
\index{método de Euler}Para $y'=f(x,y)$, $y(x_0)=y_0$ e passo $h$:
\[
x_{n+1}=x_n+h,\qquad y_{n+1}=y_n+h\,f(x_n,y_n).
\]
\end{regra}

\begin{exemplo}{1.4}{Euler na prática}
$y'=x+y$, $y(0)=1$, com $h=0{,}1$. Os passos dão $y_1=1{,}1$; $y_2=1{,}22$; $y_3=1{,}362$;
$y_4=1{,}5282$; $y_5=1{,}72102\approx y(0{,}5)$. A solução exata (seção~1.5) é
$y=2e^x-x-1$, com $y(0{,}5)\approx1{,}79744$: erro de $0{,}076$. Como a solução é convexa, as
tangentes ficam abaixo da curva e Euler \emph{subestima}. Reduzir $h$ à metade reduz o erro
aproximadamente à metade: o método tem erro proporcional a $h$.
\end{exemplo}

\figtikz{\begin{tikzpicture}
\begin{axis}[painel,width=7.4cm,height=5cm,xmin=0,xmax=0.55,ymin=0.9,ymax=1.85,xtick={0.1,0.2,0.3,0.4,0.5},ytick={1,1.5},
  xticklabel style={/pgf/number format/use comma},yticklabel style={/pgf/number format/use comma}]
\addplot[azul,thick,domain=0:0.55]{2*exp(x)-x-1};
\addplot[laranja,thick,mark=*,mark size=1.4pt] coordinates {(0,1) (0.1,1.1) (0.2,1.22) (0.3,1.362) (0.4,1.5282) (0.5,1.72102)};
\end{axis}\end{tikzpicture}}{Método de Euler para $y'=x+y$, $y(0)=1$, $h=0{,}1$: a poligonal
(verde) fica abaixo da solução exata (grafite).}

% ---------------------------------------------------------------------
\section{Equações separáveis}

\begin{definicao}
Uma equação de primeira ordem é \textbf{separável}\index{equação separável} se pode ser
escrita como $y'=g(x)\,h(y)$.
\end{definicao}

\begin{estrategia}{separação de variáveis}
\begin{enumerate}[label=\arabic*.]
\item Encontre os zeros de $h(y)$: cada $h(c)=0$ dá a solução constante $y\equiv c$.
\item Para $h(y)\neq0$, separe: $\dfrac{dy}{h(y)}=g(x)\,dx$.
\item Integre os dois lados, $\displaystyle\int\frac{dy}{h(y)}=\int g(x)\,dx$, e reúna as
  constantes numa só.
\item Se possível, isole $y$; use a condição inicial para achar a constante.
\end{enumerate}
\end{estrategia}

\begin{exemplo}{1.5}{Separando variáveis}
(a) $y'=2xy$, $y(0)=3$: $\int\frac{dy}y=\int2x\,dx$ dá $\ln|y|=x^2+C$, isto é,
$y=Ae^{x^2}$. Com $y(0)=3$: $y=3e^{x^2}$.

(b) $y'=\dfrac{x^2}{y}$: $\int y\,dy=\int x^2dx$ dá $\frac{y^2}2=\frac{x^3}3+C$, ou
$y=\pm\sqrt{\tfrac23x^3+C_1}$ (o sinal é escolhido pela condição inicial).
\end{exemplo}

\textbf{Problemas de mistura.} Um tanque contém $400$ L de água com $20$ kg de sal. Entra
salmoura com $0{,}1$ kg/L a $4$ L/min, e a mistura (homogênea) sai a $4$ L/min. Se $u(t)$ é a
massa de sal, então $u'=$ (entrada) $-$ (saída) $=0{,}4-\frac{4u}{400}=0{,}4-\frac u{100}$.
Separando, $\int\frac{du}{40-u}=\int\frac{dt}{100}$, e com $u(0)=20$:
\[
u(t)=40-20e^{-t/100}.
\]
O sal tende a $40$ kg, a concentração da salmoura que entra vezes o volume.

\textbf{Lei de resfriamento de Newton.} \index{resfriamento de Newton}Um corpo a temperatura
$T$ num ambiente a $T_a$ satisfaz $T'=k(T-T_a)$, com $k<0$. É separável:
$T(t)=T_a+(T_0-T_a)e^{kt}$.

% ---------------------------------------------------------------------
\section{A equação logística}

O modelo $P'=rP$ prevê crescimento exponencial ilimitado. Num ambiente com recursos finitos, a
população se estabiliza numa \textbf{capacidade de suporte}\index{capacidade de suporte} $K$.
A \textbf{equação logística}\index{equação logística} corrige o modelo multiplicando a taxa
por um fator que diminui à medida que $P$ se aproxima de $K$:
\[
\frac{dP}{dt}=rP\Big(1-\frac PK\Big).
\]
Os equilíbrios são $P=0$ (instável) e $P=K$ (estável). A equação é separável; com frações
parciais, $\frac{K}{P(K-P)}=\frac1P+\frac1{K-P}$, e a solução com $P(0)=P_0$ é
\[
P(t)=\frac{K\,P_0\,e^{rt}}{K-P_0+P_0\,e^{rt}} .
\]
Se $0<P_0<K$, a curva tem forma de S: cresce quase exponencialmente no início, tem
\emph{inflexão} em $P=\frac K2$ (onde o crescimento é mais rápido) e tende a $K$.

\figtikz{\begin{tikzpicture}
\begin{axis}[painel,width=8.4cm,height=5.4cm,xmin=0,xmax=30,ymin=0,ymax=5600,xtick={10,20},ytick={2000,4000},
  yticklabels={$2000$,$K=4000$},xlabel={$t$}]
\addplot[cinza,dashed,domain=0:30]{4000};
\foreach \p in {300,900,2000}{\edef\tmp{\noexpand\addplot[laranja,thick,domain=0:30,samples=100]{4000*\p*exp(0.2311*x)/(4000-\p+\p*exp(0.2311*x))};}\tmp}
\addplot[azul,thick,domain=0:30,samples=100]{4000*5500*exp(0.2311*x)/(4000-5500+5500*exp(0.2311*x))};
\end{axis}\end{tikzpicture}}{Soluções da equação logística com $r=0{,}2311$ e $K=4000$: as que
começam abaixo de $K$ (verde) crescem em S; a que começa acima (grafite) decresce até $K$.}

\begin{exemplo}{1.6}{Uma população de cervos}
Uma reserva tem $P_0=900$ cervos, taxa $r=0{,}2311$ por ano e capacidade $K=4000$. Então
$P(5)=\frac{4000\cdot900\,e^{1{,}1555}}{3100+900\,e^{1{,}1555}}\approx1919$ e
$P(10)\approx2982$. O crescimento é máximo quando $P=2000$, o que ocorre em
$t=\frac1r\ln\frac{K-P_0}{P_0}\approx5{,}35$ anos.
\end{exemplo}

% ---------------------------------------------------------------------
\section{Equações lineares de primeira ordem}

\begin{definicao}
Uma equação de primeira ordem é \textbf{linear}\index{equação linear de primeira ordem} se
pode ser escrita na forma padrão $y'+p(x)\,y=q(x)$.
\end{definicao}

A ideia é multiplicar tudo por uma função $\mu(x)$ que transforme o lado esquerdo na
derivada de um produto: $\mu y'+\mu p\,y=(\mu y)'$ exige $\mu'=\mu p$, isto é,
$\mu(x)=e^{\int p(x)\,dx}$, o \textbf{fator integrante}\index{fator integrante}.

\begin{estrategia}{equações lineares}
\begin{enumerate}[label=\arabic*.]
\item Escreva na forma padrão $y'+p(x)y=q(x)$ (divida pelo coeficiente de $y'$).
\item Calcule $\mu(x)=e^{\int p(x)\,dx}$ (sem constante).
\item Multiplique: $(\mu y)'=\mu q$. Integre: $\mu y=\int\mu q\,dx+C$.
\item Divida por $\mu$ e use a condição inicial, se houver.
\end{enumerate}
\end{estrategia}

\begin{exemplo}{1.7}{Fator integrante}
(a) $xy'+3y=4x^2-3x$, $x>0$. Forma padrão: $y'+\frac3xy=4x-3$; $\mu=e^{3\ln x}=x^3$. Então
$(x^3y)'=4x^4-3x^3$, $x^3y=\frac45x^5-\frac34x^4+C$ e
$y=\frac45x^2-\frac34x+\frac C{x^3}$.

(b) $y'-y=x$ (isto é, $y'=x+y$ do Exemplo~1.4), $y(0)=1$: $\mu=e^{-x}$,
$(e^{-x}y)'=xe^{-x}$, $e^{-x}y=-(x+1)e^{-x}+C$, $y=Ce^x-x-1$; com $y(0)=1$, $C=2$.
\end{exemplo}

\begin{exemplo}{1.8}{Queda com resistência do ar}
A velocidade de um objeto em queda satisfaz $v'=9{,}8-0{,}2v$, $v(0)=0$. Na forma padrão,
$v'+0{,}2v=9{,}8$; $\mu=e^{0{,}2t}$ e $(e^{0{,}2t}v)'=9{,}8e^{0{,}2t}$, logo
$v=49+Ce^{-0{,}2t}$ e, com $v(0)=0$, $v(t)=49\big(1-e^{-0{,}2t}\big)$. A velocidade tende à
\emph{velocidade terminal} de $49$ m/s, que é também o equilíbrio da equação.
\end{exemplo}

% ---------------------------------------------------------------------
\section*{Resumo do capítulo}\markright{Resumo do capítulo}
\addcontentsline{toc}{section}{Resumo do capítulo}
\begin{itemize}
\item Equação diferencial: envolve $y$ e suas derivadas; ordem = maior derivada. Solução geral
  (com constantes) e particular (PVI fixa as constantes).
\item Campo de direções: segmentos de inclinação $f(x,y)$; equilíbrios de $y'=f(y)$ nos zeros
  de $f$, classificados pelo sinal de $f$.
\item Euler: $y_{n+1}=y_n+h\,f(x_n,y_n)$; erro proporcional a $h$.
\item Separáveis $y'=g(x)h(y)$: $\int\frac{dy}{h(y)}=\int g(x)\,dx$; não esqueça as soluções
  constantes.
\item Logística $P'=rP(1-P/K)$: $P=\frac{KP_0e^{rt}}{K-P_0+P_0e^{rt}}$, inflexão em $K/2$.
\item Lineares $y'+py=q$: fator integrante $\mu=e^{\int p}$, $(\mu y)'=\mu q$.
\end{itemize}

% Gerado por ferramentas/gerar.py; não edite à mão.
\section*{Exercícios}\markright{Exercícios}
\addcontentsline{toc}{section}{Exercícios}
Os pontos antes de cada número indicam a dificuldade: \nivel{1} fácil, \nivel{2} médio, \nivel{3} difícil e \nivel{4} desafio. Os exercícios marcados com \estrela{} têm solução completa, com animação do gráfico, no repositório do livro: \url{https://github.com/AleffAzevedo/calculo-2-pt/blob/main/exercicios/cap1/README.md}. As respostas estão no fim do capítulo.

\begin{itemize}[leftmargin=5.2em,labelsep=0.6em,itemsep=4pt]
\item[\exnum{1}{1.1}] Dê a ordem de cada equação e diga se ela é linear: (a) $y''+3y'-y=\sen x$; (b) $(y')^2+y=x$; (c) $y'''=y\,y'$.
\item[\exnum{1}{1.2}] Verifique que $y=C_1e^{2x}+C_2e^{-x}$ é solução de $y''-y'-2y=0$ para quaisquer constantes $C_1$, $C_2$.
\item[\exnum{1}{1.3}] Para quais valores de $r$ a função $y=e^{rx}$ é solução de $y''+y'-6y=0$?
\item[\exnum{1}{1.4}] Resolva o PVI $y'=x\,e^{x^2}$, $y(0)=1$.
\item[\exnum{1}{1.5}] Resolva o PVI $y''=6x$, $y(0)=1$, $y'(0)=2$.
\item[\exnum{1}{1.6}] Uma pedra é lançada para cima, do chão, a $20$ m/s ($g=9{,}8$ m/s$^2$). Escreva $v(t)$ e $s(t)$ e encontre a altura máxima.
\item[\exnum{1}{1.7}] Encontre e classifique os equilíbrios de $y'=y^2-4$.
\item[\exnum{1}{1.8}] Encontre e classifique os equilíbrios de $y'=y(y-3)^2$.
\item[\exnum{1}{1.9}] Faça dois passos do método de Euler com $h=0{,}5$ para $y'=2x-y$, $y(0)=1$, e compare com a solução exata $y=2x-2+3e^{-x}$ em $x=1$.
\item[\exnum{1}{1.10}] Resolva $y'=\dfrac yx$ ($x>0$) por separação de variáveis.
\item[\exnum{1}{1.11}] Resolva $y'=e^{x-y}$.
\item[\exnum{1}{1.12}] Resolva $y'=1+y^2$, $y(0)=0$. Em que intervalo a solução existe?
\item[\exnum{1}{1.13}] Resolva $y'=x\,y^2$, $y(0)=1$, e diga onde a solução deixa de existir.
\item[\exnum{1}{1.14}] Um corpo a $80$\textdegree C é deixado num ambiente a $20$\textdegree C, com $T'=-0{,}1(T-20)$ ($t$ em min). Encontre $T(t)$ e quando ele chega a $40$\textdegree C.
\item[\exnum{1}{1.15}] Para $P'=0{,}5P\big(1-\tfrac P{200}\big)$, dê os equilíbrios, a capacidade de suporte e o valor de $P$ em que o crescimento é mais rápido (com a taxa nesse ponto).
\item[\exnum{1}{1.16}] Com os dados do exercício anterior e $P(0)=20$, calcule $P(4)$.
\item[\exnum{1}{1.17}] Resolva $y'+2y=4$ e descreva o comportamento das soluções quando $x\to\infty$.
\item[\exnum{1}{1.18}] Resolva $y'+\dfrac yx=3x$ ($x>0$) pelo fator integrante.
\item[\exnum{1}{1.19}] Resolva o PVI $y'-3y=e^{2x}$, $y(0)=0$.
\item[\exnum{1}{1.20}] Um tanque tem $100$ L de água pura. Entra salmoura com $0{,}2$ kg/L a $5$ L/min e sai mistura a $5$ L/min. Encontre o sal $u(t)$ e quanto há após $20$ min.
\item[\exnum{2}{1.21}] \estrela\ \textbf{Um campo que espalha as soluções.} Para $y'=y-x$, esboce o campo de direções, encontre a solução geral e mostre que só uma
solução não se afasta da reta $y=x+1$. Compare $y(0)=1$ e $y(0)=1{,}1$ em $x=5$.
\item[\exnum{2}{1.22}] \estrela\ \textbf{Euler e o número $e$.} Aplique o método de Euler a $y'=y$, $y(0)=1$, em $[0,1]$ com $n$ passos de tamanho $h=\frac1n$.
Mostre que a aproximação de $y(1)$ é $\big(1+\frac1n\big)^n$ e calcule-a para $n=1,2,4,10,100$.
\item[\exnum{2}{1.23}] \estrela\ \textbf{Soluções que acabam.} Resolva $y'=-\dfrac xy$, $y(0)=2$. Em que intervalo a solução existe? O que acontece nas pontas?
\item[\exnum{2}{1.24}] \estrela\ \textbf{A hora da morte.} Um corpo é encontrado a $30$\textdegree C numa sala a $20$\textdegree C; uma hora depois está a
$28$\textdegree C. Supondo a lei de Newton e temperatura de $37$\textdegree C no momento da morte, há quanto tempo
ocorreu a morte quando o corpo foi encontrado?
\item[\exnum{2}{1.25}] \estrela\ \textbf{Resolvendo a logística.} Resolva $P'=rP\big(1-\frac PK\big)$, $P(0)=P_0$, por separação de variáveis e frações parciais.
Para $r=0{,}2311$, $K=4000$ e $P_0=900$, quando a população cresce mais rápido?
\item[\exnum{2}{1.26}] \estrela\ \textbf{Velocidade terminal.} Um paraquedista cai com $v'=9{,}8-0{,}2v$, $v(0)=0$ (SI). Encontre $v(t)$, a distância percorrida
$s(t)$ e quanto tempo leva para atingir $95\%$ da velocidade terminal.
\item[\exnum{2}{1.27}] Classifique os equilíbrios de $y'=(y-1)(y-4)$ e esboce algumas soluções.
\item[\exnum{2}{1.28}] Use Euler com $h=0{,}25$ (quatro passos) para estimar $y(1)$ em $y'=x^2-y$, $y(0)=1$. A solução exata é $y=x^2-2x+2-e^{-x}$.
\item[\exnum{2}{1.29}] Resolva $y'=3x^2(y+1)$, $y(0)=0$.
\item[\exnum{2}{1.30}] Resolva $y'=\dfrac{\cos x}{y}$, $y(0)=2$.
\item[\exnum{2}{1.31}] Resolva $\dfrac{dy}{dx}=\dfrac{1+y^2}{1+x^2}$, $y(0)=1$, e simplifique a resposta.
\item[\exnum{2}{1.32}] Resolva $y'=y\ln y$, $y(0)=e$.
\item[\exnum{2}{1.33}] Resolva $y'+2xy=x$, $y(0)=0$.
\item[\exnum{2}{1.34}] Resolva $y'+y\tg x=\sec x$, $y(0)=1$, em $|x|<\tfrac\pi2$.
\item[\exnum{2}{1.35}] Resolva $xy'-2y=x^3e^x$ ($x>0$).
\item[\exnum{2}{1.36}] Escreva $P'=0{,}02P-0{,}00004P^2$ na forma logística e identifique $r$ e $K$.
\item[\exnum{2}{1.37}] Uma população logística tem $K=1000$, $P(0)=100$ e atinge $500$ em $t=10$. Encontre $r$.
\item[\exnum{2}{1.38}] Estime $y(1)$ para $y'=y$, $y(0)=1$, com Euler de passo $h=0{,}5$ e $h=0{,}25$. Compare os erros.
\item[\exnum{2}{1.39}] Um bolo sai do forno a $180$\textdegree C para uma sala a $25$\textdegree C e está a $100$\textdegree C após $10$ min. Quando estará a $50$\textdegree C?
\item[\exnum{2}{1.40}] Verifique que $y=\dfrac1{1+Ce^{-x}}$ resolve $y'=y(1-y)$ e encontre a solução com $y(0)=\tfrac14$.
\item[\exnum{3}{1.41}] \estrela\ \textbf{Logística com colheita.} Uma população de peixes segue $P'=0{,}4P\big(1-\frac P{1000}\big)-60$ (pesca constante de 60 por ano).
Encontre e classifique os equilíbrios. O que acontece se $P(0)=150$? Qual a maior pesca sustentável?
\item[\exnum{3}{1.42}] \estrela\ \textbf{Mistura com volume variável.} Um tanque tem $100$ L de água pura. Entra salmoura com $0{,}5$ kg/L a $3$ L/min e sai a mistura a
$2$ L/min. Encontre o sal $u(t)$ e a concentração quando o volume dobrar.
\item[\exnum{3}{1.43}] \estrela\ \textbf{Uma equação de Bernoulli.} Resolva $y'+y=y^2$ com a substituição $v=\frac1y$. Descreva as soluções com $y(0)=\tfrac13$, $y(0)=1$ e $y(0)=2$.
\item[\exnum{3}{1.44}] \estrela\ \textbf{Trajetórias ortogonais.} Encontre as curvas que cortam todas as parábolas $y=Cx^2$ em ângulo reto.
\item[\exnum{3}{1.45}] \estrela\ \textbf{Esvaziando um tanque.} Pela lei de Torricelli, o nível de um tanque cilíndrico que esvazia por um furo no fundo satisfaz
$h'=-0{,}1\sqrt h$ ($h$ em m, $t$ em min). Com $h(0)=4$, encontre $h(t)$, o tempo para esvaziar e o tempo
para o nível cair à metade.
\item[\exnum{3}{1.46}] Mostre que o PVI $y'=3y^{2/3}$, $y(0)=0$ tem pelo menos duas soluções. Por que isso não contradiz a intuição do campo de direções?
\item[\exnum{3}{1.47}] Resolva $y'=y^2-1$, $y(0)=0$, usando frações parciais.
\item[\exnum{3}{1.48}] Resolva $y'=\dfrac{x^2}{1-y^2}$ (solução implícita).
\item[\exnum{3}{1.49}] Resolva $y'+y=\sen x$.
\item[\exnum{3}{1.50}] Resolva $y'-\dfrac yx=x\sen x$ ($x>0$).
\item[\exnum{3}{1.51}] Resolva $(x+y^2)\,dy=y\,dx$ tratando $x$ como função de $y$.
\item[\exnum{3}{1.52}] Aplique Euler a $y'=-2y$, $y(0)=1$, com $h=1{,}5$. O que acontece? Para quais $h$ as aproximações decaem como a solução exata?
\item[\exnum{3}{1.53}] Mostre, derivando a equação logística, que $P''=rP'\big(1-\tfrac{2P}K\big)$ e conclua que a inflexão ocorre em $P=\tfrac K2$.
\item[\exnum{3}{1.54}] O modelo de Gompertz é $P'=rP\ln\frac KP$. Resolva-o com a substituição $u=\ln P$.
\item[\exnum{3}{1.55}] Um medicamento é infundido a $2$ mg/h e eliminado a uma taxa de $10\%$ da quantidade por hora. Com $A(0)=0$, encontre $A(t)$ e quando se atinge $90\%$ do nível de equilíbrio.
\item[\exnum{3}{1.56}] Um investimento rende $5\%$ ao ano (capitalização contínua) e recebe depósitos contínuos de R\$ $1200$ por ano. Com $S(0)=0$, quanto há após $20$ anos?
\item[\exnum{3}{1.57}] Com resistência quadrática, $v'=9{,}8-0{,}098v^2$, $v(0)=0$. Encontre $v(t)$ e a velocidade terminal.
\item[\exnum{3}{1.58}] Resolva $y'=(x+y)^2$ com a substituição $u=x+y$.
\item[\exnum{3}{1.59}] Resolva a equação homogênea $y'=\dfrac{x^2+y^2}{xy}$ ($x>0$) com $v=\frac yx$.
\item[\exnum{3}{1.60}] Uma doença se espalha numa cidade de $5000$ pessoas segundo $y'=0{,}0002\,y(5000-y)$, com $y(0)=10$. Quando metade estará infectada?
\item[\exnum{4}{1.61}] Mostre que toda solução de $y'+p(x)y=q(x)$ (com $p$, $q$ contínuas) é $y=e^{-\int p}\Big(\int e^{\int p}\,q\,dx+C\Big)$ e que o PVI com $y(x_0)=y_0$ tem exatamente uma solução.
\item[\exnum{4}{1.62}] Encontre as trajetórias ortogonais da família de circunferências $x^2+y^2=2cx$.
\item[\exnum{4}{1.63}] (Clepsidra) Pela lei de Torricelli, $A(h)\,h'=-k\sqrt h$. Que forma deve ter um recipiente de revolução para que o nível desça a taxa constante?
\item[\exnum{4}{1.64}] Um cabo suspenso satisfaz $y''=\tfrac1a\sqrt{1+(y')^2}$. Com $p=y'$, resolva e mostre que $y=a\cosh\frac xa$ (com $y'(0)=0$, $y(0)=a$).
\item[\exnum{4}{1.65}] (Riccati) Verifique que $y_1=\tfrac1x$ resolve $y'=y^2-\tfrac2{x^2}$ e encontre a solução geral com $y=y_1+\tfrac1v$.
\end{itemize}
\subsection*{Respostas}\markright{Respostas}
{\small\begin{multicols}{2}\raggedright
\begin{itemize}[leftmargin=2.7em,labelsep=0.5em,itemsep=1pt]
\item[\textbf{1.1}] (a) ordem 2, linear; (b) ordem 1, não linear; (c) ordem 3, não linear.
\item[\textbf{1.2}] $4C_1e^{2x}+C_2e^{-x}-(2C_1e^{2x}-C_2e^{-x})-2(C_1e^{2x}+C_2e^{-x})=0$.
\item[\textbf{1.3}] $r^2+r-6=0$: $r=2$ ou $r=-3$.
\item[\textbf{1.4}] $y=\tfrac12e^{x^2}+\tfrac12$.
\item[\textbf{1.5}] $y=x^3+2x+1$.
\item[\textbf{1.6}] $v=20-9{,}8t$, $s=20t-4{,}9t^2$; máximo em $t\approx2{,}04$ s, $s\approx20{,}4$ m.
\item[\textbf{1.7}] $y=-2$ estável; $y=2$ instável.
\item[\textbf{1.8}] $y=0$ instável; $y=3$ semiestável.
\item[\textbf{1.9}] $y_1=0{,}5$, $y_2=0{,}75$; exato $3/e\approx1{,}104$.
\item[\textbf{1.10}] $y=Cx$.
\item[\textbf{1.11}] $e^y=e^x+C$, ou $y=\ln(e^x+C)$.
\item[\textbf{1.12}] $y=\tg x$, para $-\tfrac\pi2<x<\tfrac\pi2$.
\item[\textbf{1.13}] $y=\dfrac2{2-x^2}$; explode em $x=\pm\sqrt2$.
\item[\textbf{1.14}] $T=20+60e^{-0{,}1t}$; $t=10\ln3\approx11{,}0$ min.
\item[\textbf{1.15}] $P=0$ e $P=200=K$; $P=100$, com $P'=25$.
\item[\textbf{1.16}] $P(4)=\dfrac{4000e^2}{180+20e^2}\approx90{,}2$.
\item[\textbf{1.17}] $y=2+Ce^{-2x}\to2$.
\item[\textbf{1.18}] $\mu=x$; $y=x^2+\dfrac Cx$.
\item[\textbf{1.19}] $y=e^{3x}-e^{2x}$.
\item[\textbf{1.20}] $u=20(1-e^{-t/20})$; $u(20)=20(1-e^{-1})\approx12{,}6$ kg.
\item[\textbf{1.21}] $y=x+1+Ce^x$; só $C=0$; $y(5)=6$ contra $y(5)\approx20{,}8$.
\item[\textbf{1.22}] $2;\ 2{,}25;\ 2{,}441;\ 2{,}594;\ 2{,}705$, tendendo a $e\approx2{,}718$.
\item[\textbf{1.23}] $y=\sqrt{4-x^2}$, para $-2<x<2$; a inclinação tende a $\mp\infty$ nas pontas.
\item[\textbf{1.24}] Cerca de $2{,}38$ h, isto é, $2$ h $23$ min.
\item[\textbf{1.25}] $P=\dfrac{KP_0e^{rt}}{K-P_0+P_0e^{rt}}$; em $t=\frac1r\ln\frac{K-P_0}{P_0}\approx5{,}35$ anos, com $P=2000$.
\item[\textbf{1.26}] $v=49(1-e^{-0{,}2t})$; $s=49t-245(1-e^{-0{,}2t})$; $t=5\ln20\approx15$ s.
\item[\textbf{1.27}] $y=1$ estável; $y=4$ instável.
\item[\textbf{1.28}] $0{,}75;\ 0{,}578;\ 0{,}496;\ 0{,}513$; exato $1-e^{-1}\approx0{,}632$.
\item[\textbf{1.29}] $y=e^{x^3}-1$.
\item[\textbf{1.30}] $y=\sqrt{2\sen x+4}$.
\item[\textbf{1.31}] $\arctg y=\arctg x+\tfrac\pi4$, ou $y=\dfrac{1+x}{1-x}$.
\item[\textbf{1.32}] $y=e^{e^x}$.
\item[\textbf{1.33}] $y=\tfrac12\big(1-e^{-x^2}\big)$.
\item[\textbf{1.34}] $\mu=\sec x$; $y=\sen x+\cos x$.
\item[\textbf{1.35}] $y=x^2(e^x+C)$.
\item[\textbf{1.36}] $P'=0{,}02P(1-P/500)$: $r=0{,}02$, $K=500$.
\item[\textbf{1.37}] $r=\tfrac{\ln9}{10}\approx0{,}220$.
\item[\textbf{1.38}] $2{,}25$ e $2{,}441$; erros $0{,}47$ e $0{,}28$ (cerca da metade).
\item[\textbf{1.39}] $t\approx25{,}1$ min.
\item[\textbf{1.40}] $C=3$: $y=\dfrac1{1+3e^{-x}}$.
\item[\textbf{1.41}] $P\approx183{,}8$ (instável) e $P\approx816{,}2$ (estável); com $P(0)=150$ a população se extingue; pesca máxima $100$ por ano.
\item[\textbf{1.42}] $u=\tfrac12(100+t)-\dfrac{500\,000}{(100+t)^2}$; em $t=100$: $87{,}5$ kg, concentração $0{,}4375$ kg/L.
\item[\textbf{1.43}] $y=\dfrac1{1+Ce^x}$; tende a $0$, fica em $1$, e explode em $x=\ln2$, respectivamente.
\item[\textbf{1.44}] As elipses $x^2+2y^2=K$.
\item[\textbf{1.45}] $h=(2-0{,}05t)^2$; esvazia em $40$ min; metade do nível em $\approx11{,}7$ min.
\item[\textbf{1.46}] $y\equiv0$ e $y=x^3$; $f(y)=3y^{2/3}$ não tem derivada limitada perto de $y=0$.
\item[\textbf{1.47}] $y=\dfrac{1-e^{2x}}{1+e^{2x}}=-\tgh x$.
\item[\textbf{1.48}] $3y-y^3=x^3+C$.
\item[\textbf{1.49}] $y=\tfrac12(\sen x-\cos x)+Ce^{-x}$.
\item[\textbf{1.50}] $y=x(C-\cos x)$.
\item[\textbf{1.51}] $\dfrac{dx}{dy}-\dfrac xy=y$: $x=y^2+Cy$.
\item[\textbf{1.52}] $y_{n+1}=(1-2h)y_n=(-2)^n$ oscila e diverge; é preciso $|1-2h|<1$, isto é, $0<h<1$.
\item[\textbf{1.53}] Derive $rP-\tfrac rKP^2$ pela regra da cadeia.
\item[\textbf{1.54}] $u'=r(\ln K-u)$: $P=K\exp\big(Ce^{-rt}\big)$.
\item[\textbf{1.55}] $A=20(1-e^{-0{,}1t})$; $t=10\ln10\approx23$ h.
\item[\textbf{1.56}] $S=24\,000(e^{0{,}05t}-1)$; $S(20)\approx$ R\$ $41\,239$.
\item[\textbf{1.57}] $v=10\tgh(0{,}98t)$; terminal $10$ m/s.
\item[\textbf{1.58}] $u'=1+u^2$: $y=\tg(x+C)-x$.
\item[\textbf{1.59}] $xv'=\tfrac1v$: $y^2=2x^2\ln x+Cx^2$.
\item[\textbf{1.60}] Logística com $r=1$, $K=5000$: $t=\ln499\approx6{,}2$ dias.
\item[\textbf{1.61}] Se $y$ é solução, $(\mu y)'=\mu q$ com $\mu=e^{\int p}>0$; integre e use $y(x_0)$ para fixar $C$.
\item[\textbf{1.62}] $y'=\tfrac{y^2-x^2}{2xy}$ para a família; as ortogonais são $x^2+y^2=2ky$.
\item[\textbf{1.63}] $A=\pi x^2$ proporcional a $\sqrt h$: $x\propto h^{1/4}$, ou $h=cx^4$.
\item[\textbf{1.64}] $\int\tfrac{dp}{\sqrt{1+p^2}}=\tfrac xa$: $p=\senh\tfrac xa$, logo $y=a\cosh\tfrac xa$.
\item[\textbf{1.65}] $v'+\tfrac2xv=-1$: $y=\dfrac1x+\dfrac{3x^2}{C-x^3}$.
\end{itemize}\end{multicols}}

